\documentclass{ximera}


\title{Area and The Definite Integral}   % Use xmlatex all -s to compile when SERVE refuses to work.
\author{Kelly Stady}
\license{CC: 0}         % replace with an appropriate license, or set it in xmPreamble

\begin{document}
\begin{abstract}
Hola
\end{abstract}
\maketitle
\label{xim:AreaActivity}


\section*{Area and the Definite Integral}

The area bounded by the $x$-axis and the function $f(x)=x^2+2$ on the interval $[0,4]$ cannot be calculated using basic geometric formulas.  
To find this area, we must begin with an approximation.

\begin{center}
    \includegraphics[width=4in,height=4in]{Area3.jpg}
    \xmalt{A graph showing the area bounded by the function f and the x-axis on the interval from 0 to 4. Area = ? is written inside the 
    shaded region.}
\end{center}

To approximate the area bounded by a function $f(x)$ and the $x$-axis on a closed interval $[a,b],$ we divide the interval into $n$ subintervals
of equal width 

\[ \Delta x = \frac{b-a}{n}. \]

We then draw $n$ rectangles whose lengths are determined by evaluating the function $f$ at a specific point within each subinterval.
The total area of the region is then approximated by summing the areas of the $n$ rectangles.

For example, we can approximate the unknown area between the $x$-axis and $f(x)=x^2+1$ on $[0,4]$ by subdividing the interval into $n = 4$
subintervals of equal length

\begin{align*} 
    \Delta x &= \frac{4-0}{4} \\
             &= 1. 
\end{align*}             

Next, we draw a rectangle on each subinterval so that its length is given by $f(c_i),$ where $c_i$ is the left endpoint of each subinterval.  
Note that while this example uses left endpoints, $c_i$ can be any $x$-value within the subinterval.

\begin{center}
    \includegraphics[height=4in,width=4in]{LowSum4.jpg}
    \xmalt{A graph showing four inscribed rectangles used to approximate the area bounded by f and the x-axis on the interval from 0 to 4.}
\end{center}

\[
\text{Area} \approx \sum_{i=1}^{n} f(c_i) \Delta x
\]



\section*{Area and the Definite Integral}

The area bounded by the $x$-axis and the function $f(x)=x^2+2$ on the interval $[0,4]$ cannot be found with a geometric formula, so we will begin by
approximating the area.

% Maple Code for Graph
% display([plot(x^2 + 2, x = 0 .. 4, 0 .. 18, color = grey, filled = true, thickness = 2), plot(x^2 + 2, x = 0 .. 4, color = black, thickness = 3)], labelfont = [COURIER, 1], % axesfont = ["TIMES", "NEW ROMAN", 17])

\begin{center}
    \includegraphics[width=4in,height=4in]{Area3.jpg} 
    \xmalt{A graph showing the area bounded by the function and the x-axis on the interval from 0 to 4.  Area = ? is written inside the shaded region.}
\end{center}

%To approximate the area bounded by a function $f(x)$ and the $x$-axis on $[a,b],$ we begin by drawing $n$ rectangles of equal width $\Delta x = \textstyle \frac{b-a}{n}$
%and lengths given by $f.$  The area is then approximated by the sum of the areas of the rectangles.

The sum of the areas of the four rectangles, shown in the graph below, can be used to approximate the area.
In the area formula shown below, the length of each rectangle is given by $f(c_i),$ where $c_i$ is the left endpoint of each subinterval.  However, $c_i$ can be any $x$-value in the $i$th subinterval.

\begin{center}
    \includegraphics[height=4in,width=4in]{LowSum4.jpg}
    \xmalt{A graph showing four inscribed rectangles used to approximate the area bounded by f and the x-axis on the interval from 0 to 4.}
\end{center}

\[ \textrm{Area} \  \approx \sum_{i=1}^{n} f(c_i) \Delta x \] 


\begin{problem}
Use the rectangles in the graph shown above to approximate the area bounded by $f(x) = x^2 + 2$ and the $x$-axis on $[0, 4].$


\begin{multipleChoice}
    \choice{\begin{align*}
        \textrm{Area} &\approx \sum_{i=1}^{4} f(c_i) \Delta x \\[1em]
        & = f(1) \cdot 1 + f(2) \cdot 1 + f(3) \cdot 1 + f(4) \cdot 1 \\[1em]
        & = 3 + 6 + 11 + 18 \\[1em]
        & = 38
        \end{align*}
    }
    \choice[correct]{\begin{align*}
        \textrm{Area} &\approx \sum_{i=1}^{4} f(c_i) \Delta x \\[1em]
        & = f(0) \cdot 1 + f(1) \cdot 1 + f(2) \cdot 1 + f(3) \cdot 1 \\[1em]
        & = 2 + 3 + 6 + 11 \\[1em]
        & = 22
        \end{align*}
    }

    \choice{None of the above.}
\end{multipleChoice}

\end{problem}

\begin{problem}
Are the rectangles in the following graph inscribed or circumscribed rectangles?

\begin{center}
    \includegraphics[height=4in,width=4in]{LowSum4.jpg}
% Alt Text
\xmalt{The graph of f and four rectangles lying inside the region bounded by f and the x-axis on the interval from 0 to 4.}
\end{center}


\begin{multipleChoice}
\choice[correct]{Inscribed}
\choice{Circumscribed}
\end{multipleChoice}


\begin{feedback}[correct]
That's right!  Inscribed rectangles lie entirely within the region we are trying to find the area of, while circumscribed rectangles extend outside the region.
Below is the graph of $f$ drawn with circumscribed rectangles.

\begin{center}

\includegraphics[height=4in,width=4in]{UpSum4.jpg}

\xmalt{The graph of f and four circumscribed rectangles drawn between f and the x-axis on the interval from 0 to 4.}
\end{center}

\end{feedback}

\end{problem}


To improve our approximation, we can increase the number of rectangles from four to twelve as shown in the graph below. 


\begin{center}
\includegraphics[height=4in,width=4in]{LowSum12.jpg}
\xmalt{The graph shows twelve rectangles used to approximate the area bounded by f and the x-axis on the interval from 0 to 4.}
\end{center}

\textbf{What happens if we use even more rectangles?}

Our approximation improves as the number of rectangles increases and if we let the number of rectangles approach infinity, we get the exact area!

\begin{definition}[\textbf{Area}]


If $f$ is continuous and nonnegative on $[a,b],$ then the area of the region bounded by
the graph of $f$ and the $x$-axis on $[a,b]$ is given by the following limit.

\[ \textrm{Area} = \lim_{n \rightarrow \infty}  \sum_{i=1}^{n} f(c_i) \Delta x, \]

where $c_i$ is any $x$-value in the $i$th subinterval and $\Delta x = \displaystyle \frac{b-a}{n}.$

\end{definition}

The sum in the above definition is called a \textbf{Riemann summ}, which is a sum used as an approximation.  
The limit of the Riemann sum in the above definition is called a \textbf{definite integral} and is denoted as follows.

\[   \lim_{n \rightarrow \infty} \sum_{i=1}^{n} f(c_i) \Delta x = \int_{a}^{b} f(x) \ dx     \]


\begin{remark}
The definite integral is not always interpreted as area.  It depends on the values of $f$ on $[a,b],$ as described below.

\begin{enumerate}
\item If $f(x)$ is \textbf{nonnegative} (positive or 0), the definite integral will give the area of the region bounded by $f$
and the $x$-axis on $[a, b].$

\item If $f$ is \textbf{negative} on $[a, b],$ the definite integral gives the negative of
the area.

\item If $f$ is \textbf{both negative and positive} on $[a, b],$ the definite integral is not interpreted
as area.
\end{enumerate}
\end{remark}

\end{document}